%%%%%%%%%%%%%%%%%%%%%%%%%%%%%%%%%%%%%%%%%%%%%%%%%%%%%%%%%%%%%%%%%%%%%%%%
%%
%%   Haskell 98 Monadic I/O Definition
%%
%%%%%%%%%%%%%%%%%%%%%%%%%%%%%%%%%%%%%%%%%%%%%%%%%%%%%%%%%%%%%%%%%%%%%%%%

%**<title>The Haskell 98 Report: Basic Input/Output</title>
%*section 7
%**~header

\chapter{Basic Input/Output}
\index{basic input/output}
\label{io}

The I/O system in \Haskell{} is purely functional, yet has all of the
expressive power found in conventional programming languages.  To
achieve this, \Haskell{} uses a \mbox{$\it monad$} to integrate I/O operations into
a purely functional context.

The I/O monad used by \Haskell{} mediates between the \mbox{$\it values$} natural to
a functional language and the \mbox{$\it actions$} that characterize I/O
operations and imperative programming in general.  The order of
evaluation of expressions in \Haskell{} is constrained only by data
dependencies; an implementation has a great deal of freedom
in choosing this order.  Actions, however, must be ordered in a
well-defined manner for program execution -- and I/O in particular --
to be meaningful.  \Haskell{}'s I/O monad provides the user with a way to
specify the sequential chaining of actions, and an implementation is
obliged to preserve this order.

\index{monad}
The term \mbox{$\it monad$} comes from a branch of mathematics known as {\em
category theory}.  From the perspective of a \Haskell{} programmer,
however, it is best to think of a monad as an {\em abstract datatype}.
In the case of the I/O monad, the abstract values are the \mbox{$\it actions$}
mentioned above.  Some operations are primitive actions,
corresponding to conventional I/O operations.  Special operations
(methods in the class \mbox{\tt Monad}, see Section~\ref{monad-class})
sequentially compose actions,
corresponding to sequencing operators (such as the semicolon) in imperative
languages.

\section{Standard I/O Functions}
\label{standard-io-functions}
Although \Haskell{} provides fairly sophisticated I/O facilities, as
defined in the \mbox{\tt IO} library, it is possible to write many
\Haskell{} programs using only the few simple functions that are
exported from the Prelude, and which are described in this section.

All I/O functions defined here are character oriented.  The treatment
of the newline character will vary on different systems.  For example,
two characters of input, return and linefeed, may read as a single
newline character.  These functions cannot be used portably for binary
I/O.

In the following, recall that \mbox{\tt String} is a synonym for \mbox{\tt [Char]} (Section~\ref{characters}).

\paragraph*{Output Functions}
These functions write to the standard output device (this is normally
the user's terminal). 
\bprog
{\small\begin{verbatim}
  putChar  :: Char -> IO ()
  putStr   :: String -> IO ()
  putStrLn :: String -> IO ()  -- adds a newline
  print    :: Show a => a -> IO ()
\end{verbatim}}
\eprog
\indextt{putChar}
\indextt{putStr}
\indextt{putStrLn}
\indextt{print}
The \mbox{\tt print} function outputs a value of any printable type to the
standard output device.
Printable types are those that are instances of class \mbox{\tt Show}; \mbox{\tt print}
converts values to strings for output using the \mbox{\tt show} operation and
adds a newline.


For example, a program to print the first 20 integers and their
powers of 2 could be written as:
\bprog
{\small\begin{verbatim}
main = print ([(n, 2^n) | n <- [0..19]])
\end{verbatim}}
\eprog

\paragraph*{Input Functions}
These functions read input from the standard input device (normally
the user's terminal).
\bprog
{\small\begin{verbatim}
  getChar     :: IO Char
  getLine     :: IO String
  getContents :: IO String
  interact    :: (String -> String) -> IO ()
  readIO      :: Read a => String -> IO a
  readLn      :: Read a => IO a
\end{verbatim}}
\eprog
\indextt{interact}%
\indextt{readIO}%
\indextt{readLn}%
\indextt{getChar}%
\indextt{getLine}%
\indextt{getContents}%
The \mbox{\tt getChar} operation raises an exception (Section~\ref{io-exceptions}) on end-of-file; a 
predicate \mbox{\tt isEOFError} that identifies this exception is defined in the \mbox{\tt IO} library.
The \mbox{\tt getLine} operation raises an exception under the same circumstances as \mbox{\tt hGetLine}, 
defined the \mbox{\tt IO} library.

The \mbox{\tt getContents} operation returns all user input as a single
string, which is read lazily as it is needed.  The \mbox{\tt interact}
function takes a function of type 
\mbox{\tt String->String} as its argument.  The entire input from the standard
input device is passed to this function
as its argument, and the resulting string is output on the
standard output device.

Typically, the \mbox{\tt read} operation from class \mbox{\tt Read} is used
to convert the string to a value.  The \mbox{\tt readIO} function is similar to \mbox{\tt read}
except that it signals parse failure to the I/O monad instead of
terminating the program.  The \mbox{\tt readLn} function combines \mbox{\tt getLine} and
\mbox{\tt readIO}.

% Removed after extended debate on the Haskell mailing list.
% By default, these input functions echo to standard output.  

% Functions in the I/O library provide full control over echoing.
% Lies, all lies

The following program simply removes all non-ASCII characters from its
standard input and echoes the result on its standard output.  (The
\mbox{\tt isAscii} function is defined in a library.)
\bprog
{\small\begin{verbatim}
main = interact (filter isAscii)
\end{verbatim}}
\eprog
\paragraph*{Files}
These functions operate on files of characters.  Files are named by
strings using some implementation-specific method to resolve strings as
file names.

The \mbox{\tt writeFile} and \mbox{\tt appendFile} functions write or append the string,
their second argument, to the file, their first argument.
The \mbox{\tt readFile} function reads a file and
returns the contents of the file as a string.  The file is read
lazily, on demand, as with \mbox{\tt getContents}.
\bprog
{\small\begin{verbatim}
  type FilePath =  String
  
  writeFile  :: FilePath -> String -> IO ()
  appendFile :: FilePath -> String -> IO ()
  readFile   :: FilePath           -> IO String
\end{verbatim}}
\eprog
\indextt{writeFile}%
\indextt{readFile}%
\indextt{appendFile}%
\indexsynonym{FilePath}%
Note that \mbox{\tt writeFile} and \mbox{\tt appendFile} write a literal string
to a file.  To write a value of any printable type, as with \mbox{\tt print}, use the
\mbox{\tt show} function to convert the value to a string first.
\bprog
{\small\begin{verbatim}
main = appendFile "squares" (show [(x,x*x) | x <- [0,0.1..2]])
\end{verbatim}}
\eprog

\section{Sequencing I/O Operations}
\indextt{>>}
\indextt{>>=}
The type constructor \mbox{\tt IO} is an instance of the \mbox{\tt Monad} class.
The two monadic binding functions, methods in the \mbox{\tt Monad} class, are
used to compose a series of I/O
operations.  The \mbox{\tt >>}
function is used where the result of the first operation is
uninteresting, for example when it is \mbox{\tt ()}.  The \mbox{\tt >>=} operation
passes the result of the first operation as an argument to the second
operation.  
\bprog
{\small\begin{verbatim}
  (>>=) :: IO a -> (a -> IO b) -> IO b 
  (>>)  :: IO a -> IO b        -> IO b
\end{verbatim}}
\eprog
For example,
\bprog
{\small\begin{verbatim}
main = readFile "input-file"                       >>= \ s ->
       writeFile "output-file" (filter isAscii s)  >>
       putStr "Filtering successful\n"
\end{verbatim}}
\eprog
is similar to the previous example using \mbox{\tt interact}, but takes its input
from \mbox{\tt "input-file"} and writes its output to \mbox{\tt "output-file"}.  A message
is printed on the standard output before the program completes.

\index{do expression}
The \mbox{\tt do} notation allows programming in a more imperative syntactic
style.  A slightly more elaborate version of the previous example
would be:
\bprog
{\small\begin{verbatim}
main = do
        putStr "Input file: "
        ifile <- getLine
        putStr "Output file: "
        ofile <- getLine
        s <- readFile ifile
        writeFile ofile (filter isAscii s) 
        putStr "Filtering successful\n"
\end{verbatim}}
\eprog

The \mbox{\tt return} function is used to define the result of an I/O
operation.  For example, \mbox{\tt getLine} is defined in terms of \mbox{\tt getChar},
using \mbox{\tt return} to define the result:
\bprog
{\small\begin{verbatim}
getLine :: IO String
getLine = do c <- getChar
             if c == '\n' then return ""
                          else do s <- getLine
                                  return (c:s)
\end{verbatim}}
\eprog

\section{Exception Handling in the I/O Monad}
\index{exception handling}
\label{io-exceptions}

The I/O monad includes a simple exception handling system.  Any I/O
operation may raise an exception instead of returning a result.

Exceptions in the I/O monad are represented by values of 
type \mbox{\tt IOError}.  This is an abstract type: its constructors are hidden
from the user.  The \mbox{\tt IO} library defines functions that construct and
examine \mbox{\tt IOError} values.  The only Prelude function that creates an
\mbox{\tt IOError} value is \mbox{\tt userError}.  User error values include a string
describing the error.
\bprog
{\small\begin{verbatim}
  userError :: String -> IOError
\end{verbatim}}
\eprog
\indextt{userError}%
Exceptions are raised and caught using the following functions:
\bprog
{\small\begin{verbatim}
  ioError :: IOError -> IO a
  catch   :: IO a    -> (IOError -> IO a) -> IO a 
\end{verbatim}}
\eprog
\indextt{ioError}%
\indextt{catch}%
The \mbox{\tt ioError} function raises an exception;
the \mbox{\tt catch} function establishes a handler that receives any
exception raised in the action protected by \mbox{\tt catch}.  An exception is
caught by the most recent handler established by \mbox{\tt catch}.  These
handlers are not selective: all exceptions are caught.  Exception
propagation must be explicitly provided in a handler by re-raising any
unwanted exceptions.  For example, in
\bprog
{\small\begin{verbatim}
f = catch g (\e -> if IO.isEOFError e then return [] else ioError e)
\end{verbatim}}
\eprog
the function \mbox{\tt f} returns \mbox{\tt []} when an end-of-file exception occurs
in \mbox{\tt g}; otherwise, the exception is propagated to the next
outer handler.  The \mbox{\tt isEOFError} function is part of \mbox{\tt IO} library.

When an exception propagates outside the main program, the \Haskell{}
system prints the associated \mbox{\tt IOError} value and exits the program.

The \mbox{\tt fail} method of the \mbox{\tt IO} instance of the \mbox{\tt Monad} class (Section~\ref{monad-class}) raises a
\mbox{\tt userError}, thus:
\bprog
{\small\begin{verbatim}
  instance Monad IO where 
    ...bindings for return, (>>=), (>>)

    fail s = ioError (userError s)
\end{verbatim}}
\eprog
\indextt{fail}%
The exceptions raised by the I/O functions in the Prelude are defined
in Chapter~\ref{module:System.IO.Error}.

%**~footer




